\documentclass[11pt,a4paper]{article}
\usepackage[utf8]{inputenc}
\usepackage[ngerman]{babel}
\usepackage{amsmath,graphicx,xcolor,listings,algorithm2e}
\definecolor{marke}{HTML}{1F6FEB}
\newcommand{\R}{\mathbb{R}}
\newcommand{\hl}[1]{\textcolor{marke}{\textbf{#1}}}
\title{Lösung des 99. Aufgabenblattes}
\author{A. Foo \and B. Bar}
\date{\today}
\begin{document}
\maketitle
\tableofcontents
\section{Aufgabe 1}\label{sec:a}
Deshalb \textbf{muss die Lösung so sein} also \textit{sozusagen}. Siehe Abschnitt~\ref{sec:a} und Gleichung~\eqref{eq:e}.\footnote{Eine Fußnote.}
Nennen wir die Variable \texttt{foo}... und \hl{wichtig} -- ``Zitat'' $x \in \R$.

\subsection{Teilaufgabe a)}
\begin{itemize}
    \item Erster Punkt
    \item Zweiter Punkt
    \begin{enumerate}
      \item innen
    \end{enumerate}
\end{itemize}
\begin{description}
    \item[Erster Punkt] Beschreibung des ersten Punktes
\end{description}
\paragraph{Teilaufgabe (i)} Text nach paragraph.
\begin{tabular}{l|c|r}
    Links & Mitte & Rechts \\
    \hline
    1 & 2 & 3 \\
\end{tabular}
\begin{figure}[H]
    \centering
    \includegraphics[width=0.5\textwidth]{example-image-a.png}
    \caption{Ein Beispielbild}
\end{figure}
\begin{lstlisting}[language=SQL]
SELECT * FROM products
WHERE price > 100;
\end{lstlisting}
\begin{algorithm}
    \KwData{this text}
    \KwResult{how to write algorithm with \LaTeX2e }
    initialization\;
    \While{not at end of this document}{
        read current\;
        \eIf{understand}{
            go to next section\;
        }{
            go back\;
        }
    }
    \caption{How to write algorithms}
\end{algorithm}
$$
E = mc^2
$$
\begin{align}
    E &= mc^2 \label{eq:e}\\
    m &= \frac{E}{c^2}
\end{align}
\% \$ \_ \# \{ \} \~{}
\end{document}
